\chapter{Conclusion}

Ce mémoire a construit un modèle interne partiel du coût en capital de la non-conformité à
DORA, où la dépendance entre les cinq piliers est portée par une cascade dirigée plutôt que
par une copule, et où le capital dépend de l'état de conformité de l'entité. Trois lignes de
résultat s'en dégagent, et ce qui est acquis se dit ici aussi nettement que ce qui
ne l'est pas.

\section{Ce qui est établi}
\paragraph{Le chiffre.} Entre l'état conforme et l'état non conforme sur les quatre canaux, le
besoin de capital passe de $6\,049$ à $20\,188$~M\euro{} à l'échelle du secteur, soit un facteur
$3{,}34$. Les canaux isolés en expliquent $9\,138$~M\euro{} et leur interaction $5\,001$, soit
$35\,\%$ de l'écart : un plan de remédiation doit donc citer l'effet de fermeture d'un canal et
non son effet isolé. La fréquence est le levier qui commande : seule, elle porte le besoin de
capital à $10\,377$~M\euro{}, et toute cible au-delà de $11\,413$~M\euro{} l'exige. Ce montant
reste un besoin de capital au sens de l'ORSA, non un SCR réglementaire.

\paragraph{Le mécanisme.} Le mécanisme tient et se distingue formellement de ses concurrents : la criticité dépend de
l'ordre de propagation, propriété qu'aucune dépendance symétrique ni aucun score
additif ne peut reproduire, et la normalisation de Leontief garantit la stabilité de la
contagion, cycles compris, par construction et non par calibration heureuse. Sur les marges
que la donnée fixe, la queue de sévérité ($\hat\xi = 0{,}5954$ au seuil publié), la
fréquence d'entrée et l'accumulation par cause commune, l'ancrage empirique est solide, et il est
éprouvé hors échantillon : la loi de comptage et la forme de la queue passent le backtest, le
niveau de la charge annuelle ne le passe pas pour une raison mesurée, la dérive de l'échelle de
sévérité. Le processus de Hawkes n'est pas distinguable de la cascade à la résolution
disponible ; il est écarté par un choix de parcimonie argumenté, non par un test qui le
réfuterait. Enfin, le modèle exprime intégralement l'effet de la
conformité sur le capital là où la Formule Standard n'en exprime rien.

\paragraph{Un acquis de méthode, qu'il faut compter comme tel.} Chaque nombre publié dans ce
mémoire est confronté, section par section, aux sorties versionnées des scripts que cette section
cite, et le contrôle rend séparément ce qu'il confirme, ce qu'il exempte avec son motif et ce
qu'il ne confirme pas. Il porte aujourd'hui sur la totalité des nombres du document, sans zone
non couverte autre que celles qui sont déclarées avec leur raison. Ce n'est pas une coulisse
d'atelier : le dispositif a trouvé trois valeurs fausses qu'aucune relecture n'avait vues, toutes
justes au moment où elles avaient été écrites et devenues fausses quand une valeur amont avait
bougé. Il est décrit au \S\ref{sec:ann-harnais}, avec ce qu'il ne garantit pas, car il ne vérifie
ni qu'un nombre est au bon endroit ni qu'il veut dire ce que la phrase prétend.

% SOURCES-SCRIPTS: 28 39 43 68 92 47 89 91 85

\section{Ce qui est démontré comme non mesurable}
La direction de la contagion n'est pas identifiable sur les données publiques. Le
mémoire l'établit sur trois sources qui échouent pour des raisons distinctes, résiste à la
tentation du test dégénéré du temps inversé, et pousse jusqu'à la meilleure expérience
naturelle disponible, le choc de chaîne d'approvisionnement MOVEit, dont l'effet net
s'annule une fois un placebo retranché. Le troisième échec ferme la dernière échappatoire :
le schéma \textsc{veris}, standard de description d'incidents adopté par la profession, porte
littéralement les deux champs manquants, et leur intersection n'est renseignée que dans
\emph{un} incident sur $10\,591$. Le problème n'est donc pas de savoir quoi collecter, il est
d'en rendre le remplissage obligatoire. Cette limite n'est pas dissimulée, elle est
convertie en méthode : le capital est borné sur l'ensemble des matrices admissibles plutôt
qu'appuyé sur une matrice posée. Le socle sans contagion est ainsi défendable sans aucune hypothèse non
identifiable, et la contagion ajoute une bande explicite au-dessus de ce socle.

% SOURCES-SCRIPTS: 30 34 28

\section{Ce qui en découle pour le régulateur}
La donnée qui rendrait la contagion mesurable se déduit par complémentarité des trois échecs :
un registre d'incidents horodatant la survenance à la semaine, catégorisant par
domaine de contrôle et non par vecteur d'attaque, consolidant une cause commune
comme un événement unique, et consignant l'ensemble des piliers touchés par chaque incident,
dont la loi jointe distingue à elle seule une matrice de sa transposée. Les deux premières exigences sont déjà des champs normalisés du
schéma \textsc{veris} : la recommandation ne demande donc pas d'inventer une taxonomie, elle
demande d'en rendre le remplissage obligatoire, ce que la notification d'incident
majeur sous DORA peut imposer et qu'une base déclarative volontaire ne peut pas. Cette
spécification n'est pas un v\oe{}u : le mémoire chiffre en euros de capital le resserrement des
bornes qu'elle apporterait, et la hiérarchise par dépendance. La limite d'identifiabilité
devient une recommandation de reporting opposable, ce qui donne au travail une portée
au-delà du modèle.

\section{Ce qui en découle pour l'entité}
Le capital n'est pas seulement un chiffre à publier, c'est un coût. En le confrontant au prix
de marché du même risque, le mémoire met les trois leviers d'une entité, remédier, détenir,
transférer, sur un seul axe en euros par an. À l'échelle d'entité, le coût annuel de détention
vaut plus de deux fois la sinistralité attendue, un traité en excès de perte correctement
dimensionné réduit le coût total d'au plus deux tiers environ, borne supérieure que la capacité
du marché, les exclusions et le risque de contrepartie entament, et la remédiation déplace ce point
d'équilibre en plus d'abaisser le capital. Le niveau reste illustratif, et il l'est doublement
puisque la descente d'échelle du secteur vers l'entité passe par deux élasticités estimées.
Mais le \emph{rapport} entre coût de détention et prix de transfert résiste mieux que le
niveau, et pour une raison qui est une identité plutôt qu'une mesure : ses deux termes
sont proportionnels au \emph{même} SCR, qui se simplifie, de sorte qu'un changement d'échelle
pur le laisse inchangé. Il vaut $2{,}78$ au taux de la directive et $3{,}51$ au taux
historique. Ce qui subsiste est un effet de forme, la descente d'échelle déplaçant la
fréquence et la sévérité séparément plutôt que de redimensionner la loi de perte, et le seul
balayage disponible le borne par le haut à $+36\,\%$ sur une plage de SCR d'un facteur quatre.
C'est la lecture décisionnelle que le travail visait, et elle est plus robuste que le niveau
sans être invariante.

% SOURCES-SCRIPTS: 58 65 67

\section{Limites, sans détour}
L'élicitation structurée est préparée et non exécutée (annexe~\ref{ch:elicitation}) : la
lecture centrale de la direction repose donc sur un classeur qualitatif, non sur une agrégation
d'experts pondérée. Deux hypothèses pèsent le plus sans être testables sur la donnée : les
niveaux de propagation, posés, et l'additivité des coûts entre piliers touchés, dont
l'élasticité vaut $0{,}95$ à l'état non conforme contre $0{,}39$ à l'état conforme
(\S\ref{sec:additivite-couts}). Dans les deux cas le signe et l'ordre de l'écart résistent à
toute la plage étudiée, non son amplitude. Le
niveau absolu du SCR est illustratif et ne doit pas être cité hors de son cadrage relatif.
L'incertitude de queue domine toute lecture ponctuelle, ce qui impose de présenter le capital
en distribution. Et l'ordre de remédiation entre les deux piliers de tête, la gouvernance et
la gestion des tiers, dépend de la direction non identifiée : le mémoire sait dire que ces
deux-là dominent, pas lequel remédier d'abord sans un dire d'expert.

% SOURCES-SCRIPTS: 74

\section{Ouvertures}
Trois prolongements se dessinent : mener l'élicitation structurée pour resserrer la bande
directionnelle sur ses deux degrés de liberté restants ; approfondir le pilier des tiers sur
des données clients, là où l'open data s'effondre ; et passer de la sensibilité de
propagation à une véritable criticité au sens de l'audit, c'est-à-dire au risque résiduel
après mesures de contrôle, une fois disponible une métrique de mitigation par pilier.

% =====================================================================
